\section{Архитектуры моделей}

Все модели строятся по схеме «предобученный трансформерный энкодер + классификационная голова». Энкодер принимает последовательность токенов и возвращает контекстуализированные представления; голова предсказывает метку для каждого токена.

\subsection{Энкодеры}

\textbf{BERT}~\cite{devlin2019bert} (\texttt{bert-base-cased}) --- предобученная модель на общем корпусе (BooksCorpus + Wikipedia), 12 слоёв, 768 скрытых размерностей.

\textbf{RoBERTa}~\cite{liu2019roberta} (\texttt{roberta-base}) --- улучшенная версия BERT с более агрессивным обучением (больше данных, динамическое маскирование, без NSP), аналогичная архитектура.

\textbf{SciBERT}~\cite{beltagy2019scibert} (\texttt{allenai/scibert\_scivocab\_cased}) --- BERT, дообученный на 1.14 млн научных статей из Semantic Scholar. Имеет собственный научный словарь. Ожидаемо превосходит общедоменные модели на научных текстах.

\subsection{Классификационные головы}

\textbf{Linear.} Дропаут $\to$ линейный слой $\mathbb{R}^{H} \to \mathbb{R}^{K}$, где $H = 768$ --- размер скрытого состояния, $K = 13$ --- число меток. Обучается с кросс-энтропией.

\textbf{MLP.} Дропаут $\to$ Linear($H$, 256) $\to$ GELU $\to$ Linear(256, $K$). Добавляет нелинейность перед финальной классификацией.

\textbf{Linear + CRF.} Линейный слой формирует эмиссионные потенциалы; CRF-слой~\cite{lafferty2001crf} моделирует зависимости между соседними метками и применяет алгоритм Витерби при декодировании. Функция потерь --- отрицательное логарифмическое правдоподобие последовательности.

\textbf{Concat-4.} Конкатенация скрытых состояний последних 4 слоёв трансформера $\mathbb{R}^{4H}$, затем дропаут и линейный слой. Позволяет модели использовать признаки разного уровня абстракции.

\subsection{Взвешивание классов}

Для компенсации дисбаланса применяются веса классов:
\[
    w_c = \sqrt{\frac{N}{K \cdot n_c}},
\]
где $N$ --- общее число токенов в обучающей выборке, $n_c$ --- число токенов класса $c$. Квадратный корень смягчает взвешивание по сравнению с обратной частотой, предотвращая чрезмерное наказание частых классов. Каждая модель обучается в двух вариантах: с весами и без.

\subsection{Полный список моделей}

\begin{table}[h]
\centering
\caption{Модели, участвующие в эксперименте}
\label{tab:models}
\begin{tabular}{llll}
\hline
\textbf{Название} & \textbf{Энкодер} & \textbf{Голова} & \textbf{Веса} \\
\hline
bert\_linear\_noweight        & BERT    & Linear      & нет \\
bert\_linear                  & BERT    & Linear      & да  \\
roberta\_linear\_noweight     & RoBERTa & Linear      & нет \\
roberta\_linear               & RoBERTa & Linear      & да  \\
scibert\_linear\_noweight     & SciBERT & Linear      & нет \\
scibert\_linear               & SciBERT & Linear      & да  \\
scibert\_mlp                  & SciBERT & MLP         & нет \\
scibert\_linear\_crf\_noweight& SciBERT & Linear+CRF  & нет \\
scibert\_linear\_crf          & SciBERT & Linear+CRF  & да  \\
scibert\_concat4              & SciBERT & Concat-4    & нет \\
\hline
\end{tabular}
\end{table}
